\documentclass[11pt]{article}
\usepackage{amsmath,amssymb,geometry}
\geometry{margin=1in}
\begin{document}

\section*{Gaussian Approximation for Color Counts in a Graph with Multiple Sources}

\subsection*{Problem setup}

\begin{itemize}
\item Graph $G = (V,E)$ with a single sink node $T$ (absorbing).
\item $K$ sources $S_1,\ldots,S_K$.
\item Each source $S_k$ injects $N_k$ particles/sec (large $N_k$), of \emph{color $k$}.
\item Particles do independent random walks and are absorbed at $T$.
\item Consider node $A \neq T$.
\item Let $X_k(A)$ be the number of color-$k$ particles at $A$ in steady state.
\end{itemize}

Goal: Approximate the \emph{joint distribution} of
\[
(X_1(A), X_2(A), \ldots, X_K(A))
\]
in steady state for large $N_k$.

\subsection*{Step 1: Mean number of particles}

The expected number of color-$k$ particles at $A$ is
\[
\lambda_k(A) = N_k \cdot H(S_k, A),
\]
where $H(S_k,A)$ is the expected number of visits to $A$ by a particle starting at $S_k$ before hitting $T$.

\subsection*{Step 2: Distribution of color counts}

Particles of color $k$ move independently:

\[
X_k(A) \sim \mathrm{Poisson}(\lambda_k(A)).
\]

\subsection*{Step 3: Joint distribution}

Because colors are independent,
\[
(X_1(A), \ldots, X_K(A)) \sim \bigotimes_{k=1}^K \mathrm{Poisson}(\lambda_k(A)).
\]

\subsection*{Step 4: Large-$\lambda_k$ Gaussian Approximation}

For large $\lambda_k(A)$:
\[
X_k(A) \approx \mathcal{N}(\lambda_k(A), \lambda_k(A)).
\]

Thus jointly:
\[
(X_1(A), \ldots, X_K(A)) \approx \mathcal{N}\bigl(\mu(A), \Sigma(A)\bigr),
\]
where
\[
\mu(A) = \begin{bmatrix}\lambda_1(A) \\ \vdots \\ \lambda_K(A)\end{bmatrix}, \quad
\Sigma(A) = \mathrm{diag}(\lambda_1(A), \ldots, \lambda_K(A)).
\]

\subsection*{Step 5: Distribution over nodes}

Suppose $A$ is random with prior $P(A)$. Then:
\[
(X_1,\ldots,X_K) \mid A \sim \mathcal{N}(\mu(A), \Sigma(A)).
\]

Joint distribution:
\[
P(A, x) = P(A) \times \mathcal{N}(x \mid \mu(A), \Sigma(A)).
\]

Marginal distribution over counts:
\[
P(x) = \sum_A P(A) \, \mathcal{N}(x \mid \mu(A), \Sigma(A)).
\]

\subsection*{Step 6: Mutual Information Approximation}

The mutual information between the color-count vector and the node identity is:
\[
I\bigl((X_1,\ldots,X_K); A\bigr) = H(X_1,\ldots,X_K) - \sum_A P(A) H\bigl(X_1,\ldots,X_K \mid A\bigr).
\]

Under the Gaussian approximation:
\[
H\bigl(X_1,\ldots,X_K \mid A\bigr) = \frac{K}{2}\log(2\pi e) + \frac{1}{2}\sum_{k=1}^K \log \lambda_k(A).
\]

The marginal distribution is a \emph{Gaussian mixture}:
\[
P(x) = \sum_A P(A) \mathcal{N}(x \mid \mu(A), \Sigma(A)).
\]

Its entropy $H(X_1,\ldots,X_K)$ typically requires numerical estimation.

\subsection*{Step 7: Final Approximate Formula}

For large $N_k$:

\[
(X_1,\ldots,X_K) \mid A \approx \mathcal{N}(\mu(A), \Sigma(A))
\]

with

\[
\mu_k(A) = N_k H(S_k, A), \quad \Sigma(A) = \mathrm{diag}(\mu(A)).
\]

Then:

\[
I \approx H(\text{Gaussian mixture}) - \sum_A P(A) \frac{1}{2}\sum_{k=1}^K \log (2\pi e \lambda_k(A)).
\]

\subsection*{Interpretation}

\begin{itemize}
\item Mutual information quantifies how well the \emph{mean color profile} $\mu(A)$ distinguishes nodes.
\item If the vectors $\mu(A)$ are well separated across nodes (relative to their variances), mutual information is large.
\item If $\mu(A)$ is nearly constant across nodes, mutual information is small.
\end{itemize}

\end{document}

%%% Local Variables:
%%% mode: latex
%%% TeX-master: t
%%% End:
